../cictikz/src/cictikz/data/tex/cictikz_preamble.tex

% A fully differential switched-capacitor sample-and-hold with a gain of
% four, and the two phase clock that runs it.
%
% Inspired by Fig. 6 of Lewis and Gray, "A pipelined 5-Msample/s 9-bit
% analog-to-digital converter", JSSC 1987, and drawn rather than
% reproduced. The original carries the common-mode feedback network and
% every dummy switch; this keeps the four elements the text argues
% about - the sampling capacitor, the feedback capacitor, the reset
% switch and the two phases - because the point being made is that the
% gain is a capacitor ratio.

\begin{circuitikz}[american, thick, transform shape]
  ../cictikz/src/cictikz/data/tex/cictikz_lib.tex

  \def\xin{0}
  \def\xsw{1.1}
  \def\xcap{2.6}
  \def\xota{5.2}
  \def\xout{8.4}
  \def\ytop{2.6}
  \def\ybot{1.0}

  \draw (\xota,\ytop) \cicOta;

  % ---- input branches: sample on 4C during phi1 ----
  \foreach \y/\lbl in {\ytop/$IN_-$, \ybot/$IN_+$} {
    \draw (\xin,\y) node[ocirc]{} -- (\xsw,\y);
    \draw (\xsw,\y) to [switch] (\xsw+0.9,\y);
    \node[scale=0.8, anchor=south] at (\xsw+0.45,\y+0.18) {$\phi_1$};
    \draw (\xsw+0.9,\y) to [C, l=$4C$] (\xota-0.5,\y);
    \node[anchor=east, scale=0.9] at (\xin-0.12,\y) {\lbl};
    %- phi2 grounds the left plate, and the charge moves to C
    \fill (\xsw+0.95,\y) circle (0.075);
  }
  %- phi2 pulls each left plate to the common mode, and the charge on
  %- 4C has nowhere to go but onto C. Routed outward, where there is room.
  \def\xcm{\xsw+0.95}
  \draw (\xcm,\ytop) -- (\xcm,\ytop+0.55)
        to [switch] (\xcm,\ytop+1.55);
  \node[scale=0.8, anchor=west] at (\xcm+0.15,\ytop+1.05) {$\phi_2$};
  \draw (\xcm-0.28,\ytop+1.55) -- (\xcm+0.28,\ytop+1.55);
  \node[scale=0.75, anchor=south] at (\xcm,\ytop+1.62) {$V_{CM}$};

  \draw (\xcm,\ybot) -- (\xcm,\ybot-0.55)
        to [switch] (\xcm,\ybot-1.55);
  \node[scale=0.8, anchor=west] at (\xcm+0.15,\ybot-1.05) {$\phi_2$};
  \draw (\xcm-0.28,\ybot-1.55) -- (\xcm+0.28,\ybot-1.55);
  \node[scale=0.75, anchor=north] at (\xcm,\ybot-1.62) {$V_{CM}$};

  % ---- feedback capacitor C with its reset switch, top and bottom ----
  % routed with right angles only: up from the summing node, across, down
  \def\yfbt{4.9}
  \def\yfbb{-1.3}

  \draw (\xota-0.5,\ytop) -- (\xota-0.5,\yfbt);
  \fill (\xota-0.5,\ytop) circle (0.075);
  \draw (\xota-0.5,\yfbt) to [C, l=$C$] (\xout-0.6,\yfbt);
  \draw (\xout-0.6,\yfbt) -- (\xout-0.6,\ytop);
  \fill (\xout-0.6,\ytop) circle (0.075);
  \draw (\xota-0.5,\yfbt+0.85) to [switch] (\xout-0.6,\yfbt+0.85);
  \node[scale=0.8, anchor=south] at ({(\xota-0.5+\xout-0.6)/2},\yfbt+1.0) {$\phi_1$};
  \draw (\xota-0.5,\yfbt) -- (\xota-0.5,\yfbt+0.85);
  \draw (\xout-0.6,\yfbt) -- (\xout-0.6,\yfbt+0.85);

  \draw (\xota-0.5,\ybot) -- (\xota-0.5,\yfbb);
  \fill (\xota-0.5,\ybot) circle (0.075);
  \draw (\xota-0.5,\yfbb) to [C, l_=$C$] (\xout-0.6,\yfbb);
  \draw (\xout-0.6,\yfbb) -- (\xout-0.6,\ybot);
  \fill (\xout-0.6,\ybot) circle (0.075);
  \draw (\xota-0.5,\yfbb-0.85) to [switch] (\xout-0.6,\yfbb-0.85);
  \node[scale=0.8, anchor=north] at ({(\xota-0.5+\xout-0.6)/2},\yfbb-1.0) {$\phi_1$};
  \draw (\xota-0.5,\yfbb) -- (\xota-0.5,\yfbb-0.85);
  \draw (\xout-0.6,\yfbb) -- (\xout-0.6,\yfbb-0.85);

  % ---- outputs ----
  \draw (cicOta_outp) -- (\xout,\ytop) node[ocirc]{};
  \draw (cicOta_outn) -- (\xout,\ybot) node[ocirc]{};
  \node[anchor=west, scale=0.9] at (\xout+0.15,\ytop) {$OUT_+$};
  \node[anchor=west, scale=0.9] at (\xout+0.15,\ybot) {$OUT_-$};

  % ---- the two phases, drawn once so the labels above mean something ----
  \begin{scope}[yshift=-5.4cm, xshift=1.2cm]
    \draw[->] (-0.4,0) -- (7.6,0) node[anchor=north east, scale=0.85] {$t$};
    \foreach \dy/\lbl/\off in {1.2/$\phi_1$/0, 0/$\phi_2$/1.6} {
      \node[anchor=east, scale=0.9] at (-0.5,\dy+0.35) {\lbl};
      \draw ({-0.4+\off},\dy) -- ({0.4+\off},\dy) -- ({0.5+\off},\dy+0.7)
            -- ({2.5+\off},\dy+0.7) -- ({2.6+\off},\dy) -- ({3.6+\off},\dy)
            -- ({3.7+\off},\dy+0.7) -- ({5.7+\off},\dy+0.7)
            -- ({5.8+\off},\dy) -- ({6.4+\off},\dy);
    }
    \node[anchor=north, scale=0.85, align=center] at (3.6,-0.35)
      {never both high at once, or the charge escapes};
  \end{scope}

\end{circuitikz}
\end{document}
